\begin{abstract}
The Goal-Frontier Maximization framework's central measure, $\volL$, quantifies
    capability \emph{possession}: the potential optionality available to a
    collective. But the experiential value it is meant to protect requires
    capability \emph{exercise}: the actual realization of that optionality in
    the world. Direct measurement of exercised optionality ($\volR$) requires
    observational access to what the collective actually does, a system that in
    its extreme is a panopticon. We reject this method in favor of a
    privacy-minimal implementation.

    We first establish a \emph{need-sufficiency architecture}: a structural
    precondition on poset construction that distinguishes involuntary
    need-satisfaction (where sacrifice magnitude reveals cost-of-access) from
    voluntary want-pursuit (where sacrifice reveals preference strength).
    Below-sufficiency sacrifices are structurally involuntary and filtered by
    the free-choice assumption before the observation channel operates. With
    this precondition in place, we introduce \emph{revealed-sacrifice
    observation}: a privacy-minimal surrogate that bounds $\volR$ from below
    using only voluntary, above-sufficiency trade events. The central result is
    the Aggregate B-to-C Lower Bound theorem: for the subset of the
    capability space covered by voluntary, above-sufficiency trade events,
    the framework lower-bounds the B-to-C gap in terms of aggregated
    sacrifice magnitudes, under assumptions S0--S2 and S4(a) of the
    revealed-preference model.  This converts a previously-uncharacterized
    divergence into a lower-bounded one, with the bound tightening
    monotonically as trade coverage accumulates.  The parallel need-sufficiency
    diagnostic tracks how efficiently the collective satisfies needs and
    whether satisfaction pathways are downstream-safe.
\end{abstract}
